\documentclass{zhvt-classic}

% 使用本地源流明体（GenRyuMin2 TW）按文件路径加载，不依赖系统字体名
% 版框（内框／界栏／书口侧栏／鱼尾）改用与版心印章同一印泥棕色：
%   印章图片 yshs.png 实色 = RGB(135,68,52) = HTML 874434
\usepackage{xcolor}
\usepackage{tikz}
% sealbrown 双注册：xcolor 供 TikZ/pgf，L3 供 cls 版框
\definecolor{sealbrown}{HTML}{874434}
\ExplSyntaxOn
\color_set:nnn { sealbrown } { HTML } { 874434 }
\ExplSyntaxOff

\zhvtset{
  grid color       = sealbrown,
  main font        = GenRyuMin2TW-SB,
  main font path   = C:/Users/Administrator/texmf/fonts/opentype/public/GenRyuMin2TW/,
jiazhu font      = GenRyuMin2TW-M,
  jiazhu font path = C:/Users/Administrator/texmf/fonts/opentype/public/GenRyuMin2TW/,
  % 页码框容纳 4 字：正常页「九十七」+前置空 1 格=4 行恰好；
  % 100+ 页「一百零一」共 5 行超出，末字「一/二/三」被框裁掉（2026-10-04 实测），故放宽到 5
  shukou page number chars = 5,
  % 书口书名框高从默认 6 字减到 4 字，把省出的 2 字高拨给 chapter 框，
  % 使 6 字书名+卷名放进 9 字高的 chapter 框后不再压线（2026-10-05）
  shukou title chars = 4,
}

\title{史记}[三王世家]
\author{汉·司马迁}
\maker{闲叟漫记·重製}


% ===== 索隐述赞样式（2026-11-08 四次定稿 · 愚大人亲手改定参数）=====
% ⚠ 竖排前提：自定义字体家族必须带 RawFeature={vertical:+vert}；
%   simhei.ttf 加 +vhal 会整体加载失败静默回退，只写 +vert；simfang 同理。
\newCJKfontfamily{\shuzanfont}[Path=C:/Windows/Fonts/,RawFeature={vertical:+vert}]{simhei.ttf}
\newCJKfontfamily{\shuzanbody}[Path=C:/Windows/Fonts/,RawFeature={vertical:+vert}]{simfang.ttf}
% 题头：152x30pt 圆角 6pt 棕底白字；baseline (0,2) 使整框落入界行格子不压版框线；
% 白字节点 x 不补偿、y +8.2pt（像素级墨迹实测 dx/dy≈0，L/R 留白约 20pt）。
\newcommand{\shuzanhead}[1]{%
  \begin{tikzpicture}[baseline={(0,2)}]
    \fill[sealbrown, rounded corners=6pt]
      (-76pt,-15pt) rectangle (76pt,15pt);
    \node[text=white, inner sep=0pt, anchor=center] at (0pt,8.2pt)
      {\shuzanfont\fontsize{23}{23}\selectfont #1};
  \end{tikzpicture}%
}
% ★ 修正书口卷名位置：chapter 框起始位置改回 cls 默认值，避免全局影响。
%   只在每个卷首页通过 \juan 命令里的空命令微调该页位置（安全，不会跨页污染）。
\ExplSyntaxOn
\AtBeginDocument{
  \dim_set:Nn \l__zhvt_shukou_chapter_title_start_dim
    { \l__zhvt_shukou_title_height_dim + \l__zhvt_font_size_dim * 2 }
  \dim_set:Nn \l__zhvt_shukou_page_number_start_dim
    { \l__zhvt_shukou_chapter_title_start_dim + \l__zhvt_shukou_chapter_title_height_dim }
}
\ExplSyntaxOff

\begin{document}

% ★★ 版心（书口）排布覆写：把书名放到鱼尾之下（2026-10-04）
%
%   cls 默认版心自上而下：书名（title 静态框，在鱼尾之上）→ 鱼尾 → 篇名（chapter 动态框）→ 页码
%   本模板改为：          鱼尾（位置不动）→ 书名 → 篇名 → 页码分割线 → 页码
%
%   手法：把 cls 放在 title 静态框里的书名，改放进 chapter 动态框（鱼尾之下）；
%         title 静态框留空。书名取 \l__zhvt_title_str，即 \title{...} 的第一参数，
%         所以换书时只改 \title，不必在此重复写书名。
%
%   ⚠ 三条必须遵守的写法（2026-10-04 实测）：
%     ① 汉字必须用 \hbox{...} 包住。flowfram 会把框内容装进一个宽 14.4mm（= 一个 line_sep
%        = 书口栏宽）的纵向 minipage 并按该宽度折行，而整页靠 PDF /Rotate 90 显示、
%        框内文字本身不变向 —— 不包 \hbox 时「围炉夜话」4 字会被逐字折成 4 行，
%        在版面上就成了横排 4 字并横向溢出版框。
%     ② 竖排下**一个 \hbox 就是一整列**。书名与篇名要同处一列，必须放进同一个 \hbox；
%        分成两个 \hbox 时，后一个会被当成新的一列、甩到隔壁栏去
%        （实测「序」跑到书口栏外侧 14.4mm 的正文列上）。中间的 \quad 在框内是横向空白，
%        显示出来正好是书名与篇名之间的纵向间隔。
%     ③ ★ 内容盒要写成 **与框等宽的定宽盒** `\hbox to \l__zhvt_shukou_chapter_title_height_dim{…\hfil}`。
%        实测：chapter 框是把内容盒沿纵向**居中**摆放的，内容宽度一变，书名就跟着漂：
%        自然宽度的 \hbox（书名 38.4mm）落在 x=157.2mm，加篇名后（57.6mm）落在 x=152.4mm，
%        目录页（无篇名）与正文页的书名就差半格。做成等宽盒后，无论有没有篇名，
%        内容固定从框顶起排，实测两页均落在 x=142.8mm，不再漂。
\ExplSyntaxOn
% 把书口 chapter 框（书名+篇名）起始位置下移 1 字高，避免 4 字卷名顶到鱼尾
% ⚠ cls 的 \__zhvt_calculate_dimensions: 在 \AtBeginDocument 里执行，会重算该 dim；
%   必须把覆写挂到 begin document 之后，否则会被二次覆盖（2026-10-05 实测）。
% (旧方案已移到导言区，直接修改 chapter 框起始位置)
% 卷标题命令 #1=书口篇名(鱼尾) #2=正文卷名（含卷序，卷名内空格用 ~）
%
% ★★ 目录页两项修订（2026-10-05，愚大人指出）★★
%   ① 页码起点统一对齐：
%      原状 —— 目录项写成 \par\hspace{4em}卷名\hspace{4em}页码，页码起点 = 卷名实宽 + 4em。
%      卷名实宽随字数变：「卷一~梁惠王上」6 字 +~、「卷七~离娄上」5 字 +~，
%      差 1 个字位 → 卷七~卷十 4 条的页码整体比其余 10 条高 27.21pt（一个字），看着上下不齐。
%      改法 —— 卷名补全角空格（宽 = 1 字 = 27.215pt）补足到 **7 字符位**再写进目录，
%      14 条页码起点遂全部落在同一条竖线上（起点对齐；末字随页码字数自然参差）。
%      ⚠ 只对目录里的副本 \l_tmpc_tl 补，正文卷名仍用 \l_tmpa_tl 原样 —— 正文像素不受影响。
%   ② 卷名整体上移半字位：前导 \hspace 由 4em 改 3.5em。
%      实测原状目录块在版框内偏下（天头留白 139.6pt / 地脚留白 110.7pt），
%      上移 0.5em = 13.6pt 后上下基本均分，视觉居中。
\NewDocumentCommand { \juan } { mm } {
	\clearpage
	\tl_set:Nn \l__zhvt_chapter_title_tl { #1 }
	\tl_set:Nn \l_tmpa_tl { #2 }
	\tl_set_eq:NN \l_tmpc_tl \l_tmpa_tl
	\int_step_inline:nn { 7 - \str_count:n { #2 } }
	  { \tl_put_right:Nn \l_tmpc_tl { 　 } }
	\tl_set:Nx \l_tmpb_tl { \zhnum { page } }
	\tl_put_left:Nn \l_tmpb_tl { \hspace { 4em } }
	\tl_put_left:Nn \l_tmpb_tl { \l_tmpc_tl }
	\tl_put_left:Nn \l_tmpb_tl { \par \hspace { 3.5em } }
	\addtocontents { toc } { \l_tmpb_tl }
	\skip_horizontal:n { \l__zhvt_font_size_dim * 4 }
	\l_tmpa_tl
	\par
	\thispagestyle { empty }
}
\cs_set_protected:Npn \__zhvt_set_shukou_content: {
	% ★ 鱼尾上方留白（title 静态框）嵌入闲叟钤印：
	%   整页 /Rotate 90，故此处把图用 \rotatebox{90} 转正为竖章；
	%   尺寸卡在版心条幅内（书内约 7mm 宽 × 33.8mm 高，留余量不压鱼尾）。
	\setstaticcontents * { title } {
		\vbox_to_ht:nn { \l__zhvt_baseline_skip_dim } {
			\vskip 14.2pt
			\hbox_to_wd:nn { \l__zhvt_shukou_title_height_dim } {
				\hfil
				\rotatebox [ origin = c ] { 90 } { \includegraphics [ width = 7mm , height = 33.8mm ] { yshs.png } }
				\hfil
			}
		}
	}
% ★ 版心（书口）统一改法（2026-11-16 愚大人定）：18pt＋框内居中＋书口文字同版框色
%   ① 原 27.2pt 时 7 字长篇名（如「廉颇蔺相如列传」）末字溢出 chapter 框、压住页码；
%      改 \fontsize{18pt}{18pt}，前后 \hfil 令字串在框内居中（不随篇名字数漂移）。
%   ② \vbox{\vskip 8.2pt} 定量下移，抵消缩字号后栏宽方向 8.2pt 上偏（像素实测 y 中心 543.3 = 原版）。
%      ⚠ \raisebox / \vspace* 在 flowfram 框内无效；换字号须同步换补偿：16pt→10.1、19pt→7.4。
%   ③ chapter 与 page~number 均置 \color{sealbrown}（HTML 874434），与版框／界行同色；正文仍黑。
%   ⚠ 页码字位、钤印、鱼尾、界行、正文一律未动；五部合订本同此口径。
	\setdynamiccontents * { chapter }
	{ \color{sealbrown} \vbox { \vskip 8.2pt \hbox to \l__zhvt_shukou_chapter_title_height_dim
		{ \hfil \fontsize{18pt}{18pt}\selectfont \l__zhvt_title_str \quad \l__zhvt_chapter_title_tl \hfil } } }
	\setdynamiccontents * { page~number } { \color{sealbrown} \quad \zhnum { page } }
}
\__zhvt_set_shukou_content:
\ExplSyntaxOff

\maketitle

% 单篇，无多卷目录，不渲染目录页

\juan{三王世家}{卷六十~三王世家第三十}
\hspace*{2em}﹁大司马臣去病昧死再拜上疏皇帝陛下。陛下过听。使臣去病待罪行间。宜专边塞之思虑。暴骸中野无以报。乃敢惟他议以干用事者。诚见陛下忧劳天下。哀怜百姓以自忘。亏膳贬乐。损郎员。皇子赖天。能胜衣趋拜。至今无号位师傅官。陛下恭让不恤。群臣私望。不敢越职而言。臣窃不胜犬马心。昧死原陛下诏有司。因盛夏吉时定皇子位。唯陛下幸察。臣去病昧死再拜以闻皇帝陛下。﹂三月乙亥。御史臣光守尚书令奏未央宫。制曰。﹁下御史。﹂\\
\hspace*{2em}六年三月戊申朔。乙亥。御史臣光守尚书令。丞非。下御史书到。言。﹁丞相臣青翟。御史大夫臣汤。太常臣充。大行令臣息。太子少傅臣安行宗正事昧死上言。大司马去病上疏曰。﹃陛下过听。使臣去病待罪行间。宜专边塞之思虑。暴骸中野无以报。乃敢惟他议以干用事者。诚见陛下忧劳天下。哀怜百姓以自忘。亏膳贬乐。损郎员。皇子赖天。能胜衣趋拜。至今无号位师傅官。陛下恭让不恤。群臣私望。不敢越职而言。臣窃不胜犬马心。昧死原陛下诏有司。因盛夏吉时定皇子位。唯原陛下幸察。﹄制曰﹃下御史﹄。臣谨与中二千石。二千石臣贺等议。古者裂地立国。并建诸侯以承天子。所以尊宗庙重社稷也。今臣去病上疏。不忘其职。因以宣恩。乃道天子卑让自贬以劳天下。虑皇子未有号位。臣青翟。臣汤等宜奉义遵职。愚憧而不逮事。方今盛夏吉时。臣青翟。臣汤等昧死请立皇子臣闳。臣旦。臣胥为诸侯王。昧死请所立国名。﹂\\
\hspace*{2em}制曰。﹁盖闻周封八百。姬姓并列。或子。男。附庸。礼﹃支子不祭﹄。云并建诸侯所以重社稷。朕无闻焉。且天非为君生民也。朕之不德。海内未洽。乃以未教成者彊君连城。即股肱何劝。其更议以列侯家之。﹂\\
\hspace*{2em}三月丙子。奏未央宫。﹁丞相臣青翟。御史大夫臣汤昧死言。臣谨与列侯臣婴齐。中二千石二千石臣贺。谏大夫博士臣安等议曰。伏闻周封八百。姬姓并列。奉承天子。康叔以祖考显。而伯禽以周公立。咸为建国诸侯。以相傅为辅。百官奉宪。各遵其职。而国统备矣。窃以为并建诸侯所以重社稷者。四海诸侯各以其职奉贡祭。支子不得奉祭宗祖。礼也。封建使守籓国。帝王所以扶德施化。陛下奉承天统。明开圣绪。尊贤显功。兴灭继绝。续萧文终之後于酂。襃厉群臣平津侯等。昭六亲之序。明天施之属。使诸侯王封君得推私恩分子弟户邑。锡号尊建百有馀国。而家皇子为列侯。则尊卑相逾。列位失序。不可以垂统於万世。臣请立臣闳。臣旦。臣胥为诸侯王。﹂三月丙子。奏未央宫。\\
\hspace*{2em}制曰。﹁康叔亲属有十而独尊者。襃有德也。周公祭天命郊。故鲁有白牡。骍刚之牲。群公不毛。贤不肖差也。﹃高山仰之。景行乡之﹄。朕甚慕焉。所以抑未成。家以列侯可。﹂\\
\hspace*{2em}四月戊寅。奏未央宫。﹁丞相臣青翟。御史大夫臣汤昧死言。臣青翟等与列侯。吏二千石。谏大夫。博士臣庆等议。昧死奏请立皇子为诸侯王。制曰。﹃康叔亲属有十而独尊者。襃有德也。周公祭天命郊。故鲁有白牡。骍刚之牲。群公不毛。贤不肖差也。﹁高山仰之。景行乡之﹂。朕甚慕焉。所以抑未成。家以列侯可。﹄臣青翟。臣汤。博士臣将行等伏闻康叔亲属有十。武王继体。周公辅成王。其八人皆以祖考之尊建为大国。康叔之年幼。周公在三公之位。而伯禽据国於鲁。盖爵命之时。未至成人。康叔後扞禄父之难。伯禽殄淮夷之乱。昔五帝异制。周爵五等。春秋三等。皆因时而序尊卑。高皇帝拨乱世反诸正。昭至德。定海内。封建诸侯。爵位二等。皇子或在襁褓而立为诸侯王。奉承天子。为万世法则。不可易。陛下躬亲仁义。体行圣德。表里文武。显慈孝之行。广贤能之路。内襃有德。外讨彊暴。极临北海。西月氏。匈奴。西域。举国奉师。舆械之费。不赋於民。虚御府之藏以赏元戎。开禁仓以振贫穷。减戍卒之半。百蛮之君。靡不乡风。承流称意。远方殊俗。重译而朝。泽及方外。故珍兽至。嘉穀兴。天应甚彰。今诸侯支子封至诸侯王。而家皇子为列侯。臣青翟。臣汤等窃伏孰计之。皆以为尊卑失序。使天下失望。不可。臣请立臣闳。臣旦。臣胥为诸侯王。﹂四月癸未。奏未央宫。留中不下。\\
\hspace*{2em}﹁丞相臣青翟。太仆臣贺。行御史大夫事太常臣充。太子少傅臣安行宗正事昧死言。臣青翟等前奏大司马臣去病上疏言。皇子未有号位。臣谨与御史大夫臣汤。中二千石。二千石。谏大夫。博士臣庆等昧死请立皇子臣闳等为诸侯王。陛下让文武。躬自切。及皇子未教。群臣之议。儒者称其术。或誖其心。陛下固辞弗许。家皇子为列侯。臣青翟等窃与列侯臣寿成等二十七人议。皆曰以为尊卑失序。高皇帝建天下。为汉太祖。王子孙。广支辅。先帝法则弗改。所以宣至尊也。臣请令史官择吉日。具礼仪上。御史奏舆地图。他皆如前故事。﹂制曰。﹁可。﹂\\
\hspace*{2em}四月丙申。奏未央宫。﹁太仆臣贺行御史大夫事昧死言。太常臣充言卜入四月二十八日乙巳。可立诸侯王。臣昧死奏舆地图。请所立国名。礼仪别奏。臣昧死请。﹂\\
\hspace*{2em}制曰。﹁立皇子闳为齐王。旦为燕王。胥为广陵王。﹂\\
\hspace*{2em}四月丁酉。奏未央宫。六年四月戊寅朔。癸卯。御史大夫汤下丞相。丞相下中二千石。二千石下郡太守。诸侯相。丞书从事下当用者。如律令。\\
\hspace*{2em}﹁维六年四月乙巳。皇帝使御史大夫汤庙立子闳为齐王。曰。於戏。小子闳。受兹青社。朕承祖考。维稽古建尔国家。封于东土。世为汉籓辅。於戏念哉。恭朕之诏。惟命不于常。人之好德。克明显光。义之不图。俾君子怠。悉尔心。允执其中。天禄永终。厥有炋臧。乃凶于而国。害于尔躬。於戏。保国艾民。可不敬与。王其戒之。﹂\\
\hspace*{2em}右齐王策。\\
\hspace*{2em}﹁维六年四月乙巳。皇帝使御史大夫汤庙立子旦为燕王。曰。於戏。小子旦。受兹玄社。朕承祖考。维稽古。建尔国家。封于北土。世为汉籓辅。於戏。荤粥氏虐老兽心。侵犯寇盗。加以奸巧边萌。於戏。朕命将率徂征厥罪。万夫长。千夫长。三十有二君皆来。降期奔师。荤粥徙域。北州以绥。悉尔心。毋作怨。毋俷德。毋乃废备。非教士不得从徵。於戏。保国艾民。可不敬与。王其戒之。﹂\\
\hspace*{2em}右燕王策。\\
\hspace*{2em}﹁维六年四月乙巳。皇帝使御史大夫汤庙立子胥为广陵王。曰。於戏。小子胥。受兹赤社。朕承祖考。维稽古建尔国家。封于南土。世为汉籓辅。古人有言曰。﹃大江之南。五湖之间。其人轻心。杨州保疆。三代要服。不及以政。﹄於戏。悉尔心。战战兢兢。乃惠乃顺。毋侗好轶。毋迩宵人。维法维则。书云。﹃臣不作威。不作福。靡有後羞。﹄於戏。保国艾民。可不敬与。王其戒之。﹂\\
\hspace*{2em}右广陵王策。\\
\hspace*{2em}太史公曰。古人有言曰﹁爱之欲其富。亲之欲其贵﹂。故王者壃土建国。封立子弟。所以襃亲亲。序骨肉。尊先祖。贵支体。广同姓於天下也。是以形势彊而王室安。自古至今。所由来久矣。非有异也。故弗论箸也。燕齐之事。无足采者。然封立三王。天子恭让。群臣守义。文辞烂然。甚可观也。是以附之世家。\\
\hspace*{2em}褚先生曰。臣幸得以文学为侍郎。好览观太史公之列传。传中称三王世家文辞可观。求其世家终不能得。窃从长老好故事者取其封策书。编列其事而传之。令後世得观贤主之指意。\\
\hspace*{2em}盖闻孝武帝之时。同日而俱拜三子为王。封一子於齐。一子於广陵。一子於燕。各因子才力智能。及土地之刚柔。人民之轻重。为作策以申戒之。谓王。﹁世为汉籓辅。保国治民。可不敬与。王其戒之。﹂夫贤主所作。固非浅闻者所能知。非博闻彊记君子者所不能究竟其意。至其次序分绝。文字之上下。简之参差长短。皆有意。人莫之能知。谨论次其真草诏书。编于左方。令览者自通其意而解说之。\\
\hspace*{2em}王夫人者。赵人也。与卫夫人并幸武帝。而生子闳。闳且立为王时。其母病。武帝自临问之。曰。﹁子当为王。欲安所置之。﹂王夫人曰。﹁陛下在。妾又何等可言者。﹂帝曰。﹁虽然。意所欲。欲於何所王之。﹂王夫人曰。﹁原置之雒阳。﹂武帝曰。﹁雒阳有武库敖仓。天下旻冲戹。汉国之大都也。先帝以来。无子王於雒阳者。去雒阳。馀尽可。﹂王夫人不应。武帝曰。﹁关东之国无大於齐者。齐东负海而城郭大。古时独临菑中十万户。天下膏腴地莫盛於齐者矣。﹂王夫人以手击头。谢曰。﹁幸甚。﹂王夫人死而帝痛之。使使者拜之曰。﹁皇帝谨使使太中大夫明奉璧一。赐夫人为齐王太后。﹂子闳王齐。年少。无有子。立。不幸早死。国绝。为郡。天下称齐不宜王云。\\
\hspace*{2em}所谓﹁受此土﹂者。诸侯王始封者必受土於天子之社。归立之以为国社。以岁时祠之。春秋大传曰。﹁天子之国有泰社。东方青。南方赤。西方白。北方黑。上方黄。﹂故将封於东方者取青土。封於南方者取赤土。封於西方者取白土。封於北方者取黑土。封於上方者取黄土。各取其色物。裹以白茅。封以为社。此始受封於天子者也。此之为主土。主土者。立社而奉之也。﹁朕承祖考﹂。祖者先也。考者父也。﹁维稽古﹂。维者度也。念也。稽者当也。当顺古之道也。\\
\hspace*{2em}齐地多变诈。不习於礼义。故戒之曰﹁恭朕之诏。唯命不可为常。人之好德。能明显光。不图於义。使君子怠慢。悉若心。信执其中。天禄长终。有过不善。乃凶于而国。而害于若身﹂。齐王之国。左右维持以礼义。不幸中年早夭。然全身无过。如其策意。\\
\hspace*{2em}传曰﹁青采出於蓝。而质青於蓝﹂者。教使然也。远哉贤主。昭然独见。诫齐王以慎内。诫燕王以无作怨。无俷德。诫广陵王以慎外。无作威与福。\\
\hspace*{2em}夫广陵在吴越之地。其民精而轻。故诫之曰﹁江湖之间。其人轻心。杨州葆疆。三代之时。迫要使从中国俗服。不大及以政教。以意御之而已。无侗好佚。无迩宵人。维法是则。无长好佚乐驰骋弋猎淫康。而近小人。常念法度。则无羞辱矣﹂。三江。五湖有鱼盐之利。铜山之富。天下所仰。故诫之曰﹁臣不作福﹂者。勿使行财币。厚赏赐。以立声誉。为四方所归也。又曰﹁臣不作威﹂者。勿使因轻以倍义也。\\
\hspace*{2em}会孝武帝崩。孝昭帝初立。先朝广陵王胥。厚赏赐金钱财币。直三千馀万。益地百里。邑万户。\\
\hspace*{2em}会昭帝崩。宣帝初立。缘恩行义。以本始元年中。裂汉地。尽以封广陵王胥四子。一子为朝阳侯。一子为平曲侯。一子为南利侯。最爱少子弘。立以为高密王。\\
\hspace*{2em}其後胥果作威福。通楚王使者。楚王宣言曰。﹁我先元王。高帝少弟也。封三十二城。今地邑益少。我欲与广陵王共发兵云。广陵王为上。我复王楚三十二城。如元王时。﹂事发觉。公卿有司请行罚诛。天子以骨肉之故。不忍致法於胥。下诏书无治广陵王。独诛首恶楚王。传曰﹁蓬生麻中。不扶自直。白沙在泥中。与之皆黑﹂者。土地教化使之然也。其後胥复祝诅谋反。自杀。国除。\\
\hspace*{2em}燕土墝埆。北迫匈奴。其人民勇而少虑。故诫之曰﹁荤粥氏无有孝行而禽兽心。以窃盗侵犯边民。朕诏将军往征其罪。万夫长。千夫长。三十有二君皆来。降旗奔师。荤粥徙域远处。北州以安矣﹂。﹁悉若心。无作怨﹂者。勿使从俗以怨望也。﹁无俷德﹂者。勿使背德也。﹁无废备﹂者。无乏武备。常备匈奴也。﹁非教士不得从徵﹂者。言非习礼义不得在於侧也。\\
\hspace*{2em}会武帝年老长。而太子不幸薨。未有所立。而旦使来上书。请身入宿卫於长安。孝武见其书。击地。怒曰。﹁生子当置之齐鲁礼义之乡。乃置之燕赵。果有争心。不让之端见矣。﹂於是使使即斩其使者於阙下。\\
\hspace*{2em}会武帝崩。昭帝初立。旦果作怨而望大臣。自以长子当立。与齐王子刘泽等谋为叛逆。出言曰。﹁我安得弟在者。今立者乃大将军子也。﹂欲发兵。事发觉。当诛。昭帝缘恩宽忍。抑案不扬。公卿使大臣请。遣宗正与太中大夫公户满意。御史二人。偕往使燕。风喻之。到燕。各异日。更见责王。宗正者。主宗室诸刘属籍。先见王。为列陈道昭帝实武帝子状。侍御史乃复见王。责之以正法。问。﹁王欲发兵罪名明白。当坐之。汉家有正法。王犯纤介小罪过。即行法直断耳。安能宽王。﹂惊动以文法。王意益下。心恐。公户满意习於经术。最後见王。称引古今通义。国家大礼。文章尔雅。谓王曰。﹁古者天子必内有异姓大夫。所以正骨肉也。外有同姓大夫。所以正异族也。周公辅成王。诛其两弟。故治。武帝在时。尚能宽王。今昭帝始立。年幼。富於春秋。未临政。委任大臣。古者诛罚不阿亲戚。故天下治。方今大臣辅政。奉法直行。无敢所阿。恐不能宽王。王可自谨。无自令身死国灭。为天下笑。﹂於是燕王旦乃恐惧服罪。叩头谢过。大臣欲和合骨肉。难伤之以法。\\
\hspace*{2em}其後旦复与左将军上官桀等谋反。宣言曰﹁我次太子。太子不在。我当立。大臣共抑我﹂云云。大将军光辅政。与公卿大臣议曰。﹁燕王旦不改过悔正。行恶不变。﹂於是脩法直断。行罚诛。旦自杀。国除。如其策指。有司请诛旦妻子。孝昭以骨肉之亲。不忍致法。宽赦旦妻子。免为庶人。传曰﹁兰根与白芷。渐之滫中。君子不近。庶人不服﹂者。所以渐然也。\\
\hspace*{2em}宣帝初立。推恩宣德。以本始元年中尽复封燕王旦两子。一子为安定侯。立燕故太子建为广阳王。以奉燕王祭祀。\\

% ===== 索隐述赞（司马贞《史记索隐》）· 统一样式 =====
% 题头「索隐述赞曰」五字并入反白框；正文每四句一列，定宽盒 \hfil 居列中；
% \fontsize{24}{41} 列距对齐全书界行（121+41k）；\clearpage 使述赞独占一页。
\clearpage
\shuzanhead{索隐述赞曰}\linebreak
{\shuzanbody\fontsize{24}{41}\selectfont
\hbox to \hsize{\hfil 三王封系。旧史烂然。褚氏後补。册书存焉。\hfil}\linebreak
\hbox to \hsize{\hfil 去病建议。青翟上言。天子冲挹。志在急贤。\hfil}\linebreak
\hbox to \hsize{\hfil 太常具礼。请立齐燕。闳国负海。旦社惟玄。\hfil}\linebreak
\hbox to \hsize{\hfil 宵人不迩。荤粥远边。明哉监戒。式防厥愆。\hfil}%
\par}
\end{document}
